% Proposed text: align repair operators with actual implementation before publication.
\subsection{Verification and Repair}
\label{sec:verification-repair}
Step~4 checks the extraction specification and its executed event log against the fixed Process Information, EEG evidence, and the source database. It preserves the selected case notion, process focus, and activity scope; revising these interpretation decisions belongs to Step~2.

Verification combines structural and execution checks with process diagnostics. The former examine specification references, timestamp roles, occurrence identity, attribution paths, SQL execution, and required output fields. The latter summarize activity coverage, repeated occurrences, duplicate identities, events per case, singleton and empty-case rates, temporal ties, and directly-follows relations. These summaries guide investigation rather than establish errors: an atomic-record case can legitimately contain one event, an unfinished occurrence can lack a completion timestamp, and independent work attributed to a shared resource can produce a complex DFG.

For a suspected missing or extra occurrence, the verifier traces the affected extraction decision to source keys, predicates, timestamp columns, and case-attribution evidence. It distinguishes a relevant recorded occurrence lost during extraction from an unrecorded occurrence, an unrelated record, or an ambiguity in interpretation. An evidence-supported repair modifies an admissible extraction decision while preserving the fixed interpretation and semantic activity labels. The revised specification is recompiled and executed, and its constraints and occurrence-level evidence are revalidated. Unsupported changes are rejected; unresolved cases remain explicit. The bounded procedure records the evidence, proposed change, and acceptance decision, rather than repairing a log solely to simplify its process model.
